\chapter{A Passion}

\lettrine{T}{he} day of his arrival, on returning from the Palais Royal, Athos, as we have seen, went straight to his hotel in the Rue Saint-Honoré. He there found the Vicomte de Bragelonne waiting for him in his chamber, chatting with Grimaud. It was not an easy thing to talk with this old servant. Two men only possessed the secret, Athos and D'Artagnan. The first succeeded, because Grimaud sought to make him speak himself; D'Artagnan, on the contrary, because he knew how to make Grimaud talk. Raoul was occupied in making him describe the voyage to England, and Grimaud had related it in all its details, with a limited number of gestures and eight words, neither more nor less. He had, at first, indicated by an undulating movement of his hand, that his master and he had crossed the sea. <Upon some expedition?> Raoul had asked.

Grimaud by bending down his head had answered, <Yes.>

<When monsieur le comte incurred much danger?> asked Raoul.

<Neither too much nor too little,> was replied by a shrug of the shoulders.

<But still, what sort of danger?> insisted Raoul.

Grimaud pointed to the sword; he pointed to the fire and to a musket that was hanging on the wall.

<Monsieur le comte had an enemy there, then?> cried Raoul.

<Monk,> replied Grimaud.

<It is strange,> continued Raoul, <that monsieur le comte persists in considering me a novice, and not allowing me to partake the honour and danger of his adventure.>

Grimaud smiled. It was at this moment Athos came in. The host was lighting him up the stairs, and Grimaud, recognizing the step of his master, hastened to meet him, which cut short the conversation. But Raoul was launched on the sea of interrogatories, and did not stop. Taking both hands of the comte, with warm, but respectful tenderness,—<How is it, monsieur,> said he, <that you have set out upon a dangerous voyage without bidding me adieu, without commanding the aid of my sword, of myself, who ought to be your support, now I have the strength; whom you have brought up like a man? Ah! monsieur, can you expose me to the cruel trial of never seeing you again?>

<Who told you, Raoul,> said the comte, placing his cloak and hat in the hands of Grimaud, who had unbuckled his sword, <who told you that my voyage was a dangerous one?>

<I,> said Grimaud.

<And why did you do so?> said Athos, sternly.

Grimaud was embarrassed; Raoul came to his assistance, by answering for him. <It is natural, monsieur, that our good Grimaud should tell me the truth in what concerns you. By whom should you be loved an supported, if not by me?>

Athos did not reply. He made a friendly motion to Grimaud, which sent him out of the room; he then seated himself in a fauteuil, whilst Raoul remained standing before him.

<But it is true,> continued Raoul, <that your voyage was an expedition, and that steel and fire threatened you?>

<Say no more about that, vicomte,> said Athos, mildly. <I set out hastily, it is true: but the service of King \regnal{Charles II} required a prompt departure. As to your anxiety, I thank you for it, and I know that I can depend on you. You have not wanted for anything, vicomte, in my absence, have you?>

<No, monsieur, thank you.>

<I left orders with Blaisois to pay you a hundred pistoles, if you should stand in need of money.>

<Monsieur, I have not seen Blaisois.>

<You have been without money, then?>

<Monsieur, I had thirty pistoles left from the sale of the horses I took in my last campaign, and M.~le Prince had the kindness to allow me to win two hundred pistoles at his play-table three months ago.>

<Do you play? I don't like that, Raoul.>

<I never play, monsieur; it was M.~le Prince who ordered me to hold his cards at Chantilly—one night when a courier came to him from the king. I won, and M.~le Prince commanded me to take the stakes.>

<Is that a practice in the household, Raoul?> asked Athos with a frown.

<Yes, monsieur; every week M.~le Prince affords, upon one occasion or another, a similar advantage to one of his gentlemen. There are fifty gentlemen in his highness's household; it was my turn.>

<Very well! You went into Spain, then?>

<Yes, monsieur, I made a very delightful and interesting journey.>

<You have been back a month, have you not?>

<Yes, monsieur.>

<And in the course of that month?>

<In that month\longdash>

<What have you done?>

<My duty, monsieur.>

<Have you not been home, to La Fère?>

Raoul coloured. Athos looked at him with a fixed but tranquil expression.

<You would be wrong not to believe me,> said Raoul. <I feel that I coloured, and in spite of myself. The question you did me the honour to ask me is of a nature to raise in me much emotion. I colour, then, because I am agitated, not because I meditate a falsehood.>

<I know, Raoul, you never lie.>

<No, monsieur.>

<Besides, my young friend, you would be wrong; what I wanted to say\longdash>

<I know quite well, monsieur. You would ask me if I have not been to Blois?>

<Exactly so.>

<I have not been there; I have not even seen the person to whom you allude.>

Raoul's voice trembled as he pronounced these words. Athos, a sovereign judge in all matters of delicacy, immediately added, <Raoul, you answer me with a painful feeling; you are unhappy.>

<Very, monsieur; you have forbidden me to go to Blois, or to see Mademoiselle de La Vallière again.> Here the young man stopped. That dear name, so delightful to pronounce, made his heart bleed, although so sweet upon his lips.

<And I have acted rightly, Raoul.> Athos hastened to reply. <I am neither an unjust nor a barbarous father; I respect true love; but I look forward for you to a future—an immense future. A new reign is about to break upon us like a fresh dawn. War calls upon a young king full of chivalric spirit. What is wanting to assist this heroic ardour is a battalion of young and free lieutenants who would rush to the fight with enthusiasm, and fall, crying: <Vive le Roi!> instead of <Adieu, my dear wife.> You understand that, Raoul. However brutal my reasoning may appear, I conjure you, then, to believe me, and to turn away your thoughts from those early days of youth in which you took up this habit of love—days of effeminate carelessness, which soften the heart and render it incapable of consuming those strong bitter draughts called glory and adversity. Therefore, Raoul, I repeat to you, you should see in my counsel only the desire of being useful to you, only the ambition of seeing you prosper. I believe you capable of becoming a remarkable man. March alone, and you will march better, and more quickly.>

<You have commanded, monsieur,> replied Raoul, <and I obey.>

<Commanded!> cried Athos. <Is it thus you reply to me? I have commanded you! Oh! you distort my words as you misconceive my intentions. I do not command you; I request you.>

<No, monsieur, you have commanded,> said Raoul, persistently; <had you requested me, your request is even more effective than your order. I have not seen Mademoiselle de La Vallière again.>

<But you are unhappy! you are unhappy!> insisted Athos.

Raoul made no reply.

<I find you pale; I find you dull. The sentiment is strong, then?>

<It is a passion,> replied Raoul.

<No—a habit.>

<Monsieur, you know I have travelled much, that I have passed two years far away from her. A habit would yield to an absence of two years, I believe; whereas, on my return, I loved not more, that was impossible, but as much. Mademoiselle de La Vallière is for me the one lady above all others; but you are for me a god upon earth—to you I sacrifice everything.>

<You are wrong,> said Athos; <I have no longer any right over you. Age has emancipated you; you no longer even stand in need of my consent. Besides, I will not refuse my consent after what you have told me. Marry Mademoiselle de La Vallière, if you like.>

Raoul was startled, but suddenly: <You are very kind, monsieur,> said he; <and your concession excites my warmest gratitude, but I will not accept it.>

<Then you now refuse?>

<Yes, monsieur.>

<I will not oppose you in anything, Raoul.>

<But you have at the bottom of your heart an idea against this marriage: it is not your choice.>

<That is true.>

<That is sufficient to make me resist: I will wait.>

<Beware, Raoul! What you are now saying is serious.>

<I know it is, monsieur; as I said, I will wait.>

<Until I die?> said Athos, much agitated.

<Oh! monsieur,> cried Raoul, with tears in his eyes, <is it possible that you should wound my heart thus? I have never given you cause of complaint!>

<Dear boy, that is true,> murmured Athos, pressing his lips violently together to conceal the emotion of which he was no longer master. <No, I will no longer afflict you; only I do not comprehend what you mean by waiting. Will you wait till you love no longer?>

<Ah! for that!—no, monsieur. I will wait till you change your opinion.>

<I should wish to put the matter to a test, Raoul; I should like to see if Mademoiselle de La Vallière will wait as you do.>

<I hope so, monsieur.>

<But, take care, Raoul! suppose she did not wait? Ah, you are young, so confiding, so loyal! Women are changeable.>

<You have never spoken ill to me of women, monsieur; you have never had to complain of them; why should you doubt of Mademoiselle de La Vallière?>

<That is true,> said Athos, casting down his eyes; <I have never spoken ill to you of women; I have never had to complain of them; Mademoiselle de La Vallière never gave birth to a suspicion; but when we are looking forward, we must go even to exceptions, even to improbabilities! If, I say, Mademoiselle de La Vallière should not wait for you?>

<How, monsieur?>

<If she turned her eyes another way.>

<If she looked favourably upon another, do you mean, monsieur?> said Raoul, pale with agony.

<Exactly.>

<Well, monsieur, I would kill him,> said Raoul, simply, <and all the men whom Mademoiselle de La Vallière should choose, until one of them had killed me, or Mademoiselle de La Vallière had restored me her heart.>

Athos started. <I thought,> resumed he, in an agitated voice, <that you called my just now your god, your law in this world.>

<Oh!> said Raoul, trembling, <you would forbid me the duel?>

<Suppose I did forbid it, Raoul?>

<You would not forbid me to hope, monsieur; consequently you would not forbid me to die.>

Athos raised his eyes toward the vicomte. He had pronounced these words with the most melancholy look. <Enough,> said Athos, after a long silence, <enough of this subject, upon which we both go too far. Live as well as you are able, Raoul, perform your duties, love Mademoiselle de La Vallière; in a word, act like a man, since you have attained the age of a man; only do not forget that I love you tenderly, and that you profess to love me.>

<Ah! monsieur le comte!> cried Raoul, pressing the hand of Athos to his heart.

<Enough, dear boy, leave me; I want rest. À propos, M.~d'Artagnan has returned from England with me; you owe him a visit.>

<I will pay it, monsieur, with great pleasure. I love Monsieur d'Artagnan exceedingly.>

<You are right in doing so; he is a worthy man and a brave cavalier.>

<Who loves you dearly.>

<I am sure of that. Do you know his address?>

<At the Louvre, I suppose, or wherever the king is. Does he not command the musketeers?>

<No; at present M.~d'Artagnan is absent on leave; he is resting for awhile. Do not, therefore, seek him at the posts of his service. You will hear of him at the house of a certain Planchet.>

<His former lackey?>

<Exactly; turned grocer.>

<I know; Rue des Lombards?>

<Somewhere thereabouts, or Rue des Arcis.>

<I will find it, monsieur—I will find it.>

<You will say a thousand kind things to him, on my part, and ask him to come and dine with me before I set out for La Fère.>

<Yes, monsieur.>

<Good-might, Raoul!>

<Monsieur, I see you wear an order I never saw you wear before; accept my compliments.>

<The Fleece!—that is true. A bauble, my boy, which no longer amuses an old child like myself. Good-night, Raoul!> 